\begin{argument}
\premises{
  % -----------------------------
  % NLP and skills extraction
  % -----------------------------
  \item There exist validated NLP methods for extracting technical skills from job postings with high precision and recall \cite{zhang-etal-2022-skillspan}.
  \item Contextual embeddings (e.g., Sentence-BERT) map job postings and repository descriptions into a shared feature space where cosine similarity reflects semantic similarity, as validated on semantic textual similarity benchmarks \cite{reimers_2019_SentenceBERTSentence}.
  \item Technical skills and repository descriptions are expressed in relatively standardized terminology, enabling reliable extraction and semantic comparison \cite{esco2023taxonomy}.

  % -----------------------------
  % GitHub activity as a signal
  % -----------------------------
  \item Stars indicate repository visibility and perceived popularity, while forks indicate community engagement and experimentation with the repository \cite{dabbish2012social}.
  \item Sustained repository activity (commits, issues, pull requests) signals ongoing experimentation and adoption of new technologies \cite{borges_2016_UnderstandingFactors}.
  \item Early engagement with repositories may reflect emerging technology trends that are later reflected in labor market demand.

  % -----------------------------
  % Job postings as demand proxies
  % -----------------------------
  \item Job postings represent explicit employer demand for skills and thus serve as a proxy for labor market demand trends.
  \item Temporal trends in job postings correlate with broader industry hiring trends \cite{radlinski2011learning,zhang-etal-2022-skillspan}.
  \item Skill mentions in job postings are described with relatively consistent terminology, enabling robust extraction using NLP techniques \cite{esco2023taxonomy}.

  % -----------------------------
  % Embeddings & similarity
  % -----------------------------
  \item Cosine similarity between embeddings in the same feature space provides a quantitative measure of semantic similarity suitable for ranking items \cite{mikolov2013distributed, reimers_2019_SentenceBERTSentence}.
  \item Cosine similarity preserves approximate local neighborhoods: if two items are both highly similar to a shared reference embedding, they are likely to be semantically related to each other \cite{ethayarajh2019contextual}.
  \item Cosine similarity induces a partial ordering of job postings with respect to each repository: if $\cos(J_A, R) > \cos(J_B, R)$, then $J_A$ is considered more semantically aligned with $R$ than $J_B$, enabling temporal aggregation and trend detection.
  \item Aggregated similarity between repositories and job postings can be used to detect alignment between emerging technologies and labor market demand.
  \item Increases in repository activity may precede increases in job posting similarity, justifying temporal alignment and forecasting analysis.
}

\end{argument}